\chapter{3 次元化：球面座標の肩と SO(3)（Sim07）}\label{ch:3d}

\section{座標の選び方}

3 次元の角運動量と特異点の式は付録~\ref{app:3d} を参照。

\begin{figure}[htbp]\centering
\resizebox{0.6\linewidth}{!}{% 球面座標の肩と円錐運動
\tdplotsetmaincoords{65}{120}
\begin{tikzpicture}[>=Stealth,font=\small,tdplot_main_coords,scale=2.4]
\draw[->] (0,0,0)--(1.3,0,0) node[anchor=north east]{$x$（前）};
\draw[->] (0,0,0)--(0,1.2,0) node[anchor=north west]{$y$（左）};
\draw[->] (0,0,0)--(0,0,1.2) node[anchor=south]{$z$（上）};
\pgfmathsetmacro{\psA}{40}\pgfmathsetmacro{\alA}{30}
\coordinate (U) at ({cos(\alA)*cos(\psA)},{cos(\alA)*sin(\psA)},{sin(\alA)});
\coordinate (P) at ({cos(\alA)*cos(\psA)},{cos(\alA)*sin(\psA)},0);
\draw[line width=2.5pt,blue!60] (0,0,0)--(U);
\draw[dashed] (0,0,0)--(P)--(U);
\tdplotdrawarc[->]{(0,0,0)}{0.45}{0}{\psA}{anchor=north}{$\psi$}
\node at (0.62,0.55,0.25) {$\alpha$};
\draw[red,thick] plot[domain=0:360,samples=72,variable=\t] ({cos(30)},{sin(30)*cos(\t)},{sin(30)*sin(\t)});
\draw[red,->,thick] (0,0,0)--(1.25,0,0);
\node[red,font=\footnotesize,anchor=east] at (0.9,-0.6,0.5) {円錐（軸：前方）};
\fill (0,0,0) circle (0.8pt) node[anchor=north east,font=\footnotesize]{肩};
\node[blue!70!black,font=\footnotesize,anchor=south west] at (U) {腕 $\bm u(\psi,\alpha)$};
\fill[gray] (0,0,1) circle (0.6pt) node[anchor=west,font=\footnotesize]{極（特異点）};
\end{tikzpicture}
}
\caption{肩を原点とする球面座標 $(\psi,\alpha)$ と、前方を軸とする円錐運動。真上（極）は座標の特異点になる。}\label{fig:tikz-sph}
\end{figure}


\begin{longtable}[]{@{}
  >{\raggedright\arraybackslash}p{(\columnwidth - 4\tabcolsep) * \real{0.3333}}
  >{\raggedright\arraybackslash}p{(\columnwidth - 4\tabcolsep) * \real{0.3333}}
  >{\raggedright\arraybackslash}p{(\columnwidth - 4\tabcolsep) * \real{0.3333}}@{}}
\toprule\noalign{}
\begin{minipage}[b]{\linewidth}\raggedright
対象
\end{minipage} & \begin{minipage}[b]{\linewidth}\raggedright
座標
\end{minipage} & \begin{minipage}[b]{\linewidth}\raggedright
理由
\end{minipage} \\
\midrule\noalign{}
\endhead
\bottomrule\noalign{}
\endlastfoot
腕の向き & \textbf{肩を原点とする球面座標} (ψ: 方位角, α: 仰角)、u = (cos α cos ψ, cos α sin ψ, sin α) & 2 自由度の肩（旋回＋俯仰）の関節角そのもの。旋回動作を自然に書ける \\
胴体の姿勢 & \textbf{回転行列（SO(3)）} & 3 次元の回転は順序で結果が変わる（非可換）。オイラー角は特異点を持つため、積分には使わない \\
\end{longtable}

\begin{itemize}
\tightlist
\item
  \textbf{球面座標の弱点}: α = ±90°（真上・真下）で ψ が定まらない（座標の特異点）。 これは座標の都合だけでなく、2 軸の肩（ジンバル）が実際に持つ \textbf{ジンバルロック} と同じ現象である（Part A3）。
\item
  \textbf{モデルの検証}: 腕を x-z 平面内（ψ = 0）だけで動かすと、平面モデル（Sim01）と 5 桁一致した。
\item
  \textbf{腕のモデル化}: 腕（一様な棒）は 40 個の点質量で表し、角運動量と反力を点質量の和として計算した。
\end{itemize}

3 次元での AMBAC の基本式は、平面の式（第~\ref{ch:theory}章）がベクトル・行列になったものである。

\[
\mathbf L = A(q)\,\boldsymbol\omega + B(q)\,\dot q = 0 \quad\Rightarrow\quad
\boldsymbol\omega = -A(q)^{-1}B(q)\,\dot q,\qquad \dot R = R\,[\boldsymbol\omega]_\times
\]

\begin{itemize}
\tightlist
\item
  A(q): 3×3 の慣性テンソル（形状 q で固まった系全体）
\item
  B(q): 3×2 の結合行列
\item
  ω: 胴体座標で表した胴体の角速度
\end{itemize}

\section{Part A: 宇宙での 3 次元 AMBAC（人間大）}

\FIG{sim07_space.png}

\subsection{A1: 腕で円錐を描く（宇宙飛行士の「腕回し」）}

\begin{longtable}[]{@{}
  >{\raggedright\arraybackslash}p{(\columnwidth - 4\tabcolsep) * \real{0.3333}}
  >{\raggedright\arraybackslash}p{(\columnwidth - 4\tabcolsep) * \real{0.3333}}
  >{\raggedright\arraybackslash}p{(\columnwidth - 4\tabcolsep) * \real{0.3333}}@{}}
\toprule\noalign{}
\begin{minipage}[b]{\linewidth}\raggedright
半頂角 β
\end{minipage} & \begin{minipage}[b]{\linewidth}\raggedright
1 周あたりの胴体回転
\end{minipage} & \begin{minipage}[b]{\linewidth}\raggedright
回転軸
\end{minipage} \\
\midrule\noalign{}
\endhead
\bottomrule\noalign{}
\endlastfoot
10° & 0.65° & ほぼ −x（円錐の軸と逆向き） \\
30° & 5.7° & 〃 \\
45° & 12.4° & 〃 \\
70° & 25.4° & 〃 \\
\end{longtable}

\begin{itemize}
\tightlist
\item
  \textbf{2 自由度の肩 1 つ（腕 1 本）だけで、任意の角度まで回れる。} 平面では関節 2 つ（腕＋脚）が必要だったが、 3 次元の肩はそれだけで 2 自由度の形状空間を持つ。
\item
  小さい円錐では、回転角 ∝ β\^{}2.00（理論どおりの 2 次効果）。Sim05 で議論した「幾何学的位相は振幅の 2 次効果」が、3 次元でも厳密に成り立つ。
\item
  円錐を逆回りにすると、回転は正確に符号反転する。
\end{itemize}

\subsection{A2: 円錐の軸の向き（β = 45°）}

\begin{longtable}[]{@{}
  >{\raggedright\arraybackslash}p{(\columnwidth - 4\tabcolsep) * \real{0.3333}}
  >{\raggedright\arraybackslash}p{(\columnwidth - 4\tabcolsep) * \real{0.3333}}
  >{\raggedright\arraybackslash}p{(\columnwidth - 4\tabcolsep) * \real{0.3333}}@{}}
\toprule\noalign{}
\begin{minipage}[b]{\linewidth}\raggedright
円錐の軸
\end{minipage} & \begin{minipage}[b]{\linewidth}\raggedright
1 周あたり
\end{minipage} & \begin{minipage}[b]{\linewidth}\raggedright
胴体の回転軸と円錐の軸のずれ
\end{minipage} \\
\midrule\noalign{}
\endhead
\bottomrule\noalign{}
\endlastfoot
前方 & 12.4° & 17° \\
右横 & 11.1° & 20° \\
真上（方位角を 1 周させる旋回） & 58.0° & 8° \\
前上方 45° & 39.7° & 31° \\
\end{longtable}

\begin{itemize}
\tightlist
\item
  胴体はおおむね円錐の軸の逆向きに回るが、\textbf{ぴったり逆向きではない}（8〜31° ずれる）。 肩が胴体の中心からずれていること、胴体の慣性テンソルが非等方であることによる。
\item
  「真上」軸の円錐は、方位角 ψ を 360° 回し続ける動作で、腕が胴体のまわりを 1 周する。 胴体の主軸（縦軸）まわりの慣性が最も小さいので、回転量が最も大きい。
\end{itemize}

\subsection{A3: 球面座標の特異点（ジンバルロック）}

円錐が極（真上）の近くを通ると、腕の先端は一定の速さで動いているのに、方位角の関節速度 ψ' が発散する。

\begin{longtable}[]{@{}
  >{\raggedright\arraybackslash}p{(\columnwidth - 4\tabcolsep) * \real{0.3333}}
  >{\raggedright\arraybackslash}p{(\columnwidth - 4\tabcolsep) * \real{0.3333}}
  >{\raggedright\arraybackslash}p{(\columnwidth - 4\tabcolsep) * \real{0.3333}}@{}}
\toprule\noalign{}
\begin{minipage}[b]{\linewidth}\raggedright
ループと極との最小角距離
\end{minipage} & \begin{minipage}[b]{\linewidth}\raggedright
最大 ψ'
\end{minipage} & \begin{minipage}[b]{\linewidth}\raggedright
最大 α'
\end{minipage} \\
\midrule\noalign{}
\endhead
\bottomrule\noalign{}
\endlastfoot
30° & 3.1 rad/s & 1.57 rad/s \\
3° & 30 rad/s & 1.57 rad/s \\
0.3° & 300 rad/s & 1.57 rad/s \\
\end{longtable}

ψ' は極からの距離に反比例する。\textbf{球面座標は旋回動作には便利だが、極の近くでは使えない。} 対策は次のとおり。 - 計画は腕の向きのベクトル u（または四元数）で行い、関節角への変換は極から離れた領域だけで行う。 - 肩の軸の取り付け方向を工夫し、極が作業領域の外に来るようにする（実機のジンバル設計と同じ考え方）。 - 冗長な関節（3 自由度の肩など）を持たせる。

\subsection{A4: 回転の非可換性}

ヨー用の円錐 1 周（約 58°）とピッチ用の円錐 3 周（約 33°）を、順序を入れ替えて実行した。

\begin{longtable}[]{@{}
  >{\raggedright\arraybackslash}p{(\columnwidth - 2\tabcolsep) * \real{0.5000}}
  >{\raggedright\arraybackslash}p{(\columnwidth - 2\tabcolsep) * \real{0.5000}}@{}}
\toprule\noalign{}
\begin{minipage}[b]{\linewidth}\raggedright
順序
\end{minipage} & \begin{minipage}[b]{\linewidth}\raggedright
最終姿勢（回転ベクトル）
\end{minipage} \\
\midrule\noalign{}
\endhead
\bottomrule\noalign{}
\endlastfoot
ヨー → ピッチ & (8.2°, 41.4°, −64.1°) \\
ピッチ → ヨー & (−22.4°, 33.7°, −65.3°) \\
\textbf{2 つの最終姿勢の差} & \textbf{29.4°} \\
\end{longtable}

平面ではループの順序を入れ替えても結果は同じだった（回転が 1 軸だけで可換）。 \textbf{3 次元では動作の順序そのものが結果を変えるので、姿勢変更の計画は「どの動作をどの順で行うか」という系列の問題になる。} これは学習・計画（フィジカル AI）にとって、平面より本質的に難しい点である。

\section{Part B: 地上での 3 次元（戦車型台車＋球面座標の腕）}

\FIG{sim07_ground.png}

\subsection{B1: 静的な定格荷重マップ}

腕の向き（ψ, α）ごとに、静止状態で転倒しない最大荷重を求めた（安全余裕 5 cm）。 クレーンや油圧ショベルの \textbf{定格荷重表} に相当する。

\begin{longtable}[]{@{}
  >{\raggedright\arraybackslash}p{(\columnwidth - 2\tabcolsep) * \real{0.5000}}
  >{\raggedright\arraybackslash}p{(\columnwidth - 2\tabcolsep) * \real{0.5000}}@{}}
\toprule\noalign{}
\begin{minipage}[b]{\linewidth}\raggedright
腕の向き（水平）
\end{minipage} & \begin{minipage}[b]{\linewidth}\raggedright
最大荷重
\end{minipage} \\
\midrule\noalign{}
\endhead
\bottomrule\noalign{}
\endlastfoot
前（ψ = 0°） & 41.5 kg \\
斜め前（45°） & 42.2 kg \\
\textbf{横（90°）} & \textbf{24.3 kg} \\
後（180°） & 41.5 kg \\
\end{longtable}

\begin{itemize}
\tightlist
\item
  接地面が前後に長く左右に狭いので、\textbf{横向きが最も弱い}（前後の約 6 割）。
\item
  腕を高く上げる（α ≳ 70°）と荷物が接地面の真上に来て、静的には転倒しなくなる（図上部の空白）。
\end{itemize}

\subsection{B2: 腕を水平に 180° 旋回（前 → 左 → 後、荷物 10 kg）}

すべての方策に同じ制約（転倒・滑り・関節速度 6 rad/s・加速度 40 rad/s²・仰角 ≥ −35°）を課し、最短時間を比べた。

\begin{longtable}[]{@{}
  >{\raggedright\arraybackslash}p{(\columnwidth - 4\tabcolsep) * \real{0.3333}}
  >{\raggedright\arraybackslash}p{(\columnwidth - 4\tabcolsep) * \real{0.3333}}
  >{\raggedright\arraybackslash}p{(\columnwidth - 4\tabcolsep) * \real{0.3333}}@{}}
\toprule\noalign{}
\begin{minipage}[b]{\linewidth}\raggedright
方策
\end{minipage} & \begin{minipage}[b]{\linewidth}\raggedright
地球
\end{minipage} & \begin{minipage}[b]{\linewidth}\raggedright
月
\end{minipage} \\
\midrule\noalign{}
\endhead
\bottomrule\noalign{}
\endlastfoot
水平のまま旋回 & 1.83 s & 4.48 s \\
途中で 60° 持ち上げる & 1.24 s & 3.04 s \\
途中で 35° 下げる & 1.01 s & 2.48 s \\
\textbf{最適化（仰角の時間変化と旋回のタイミング）} & \textbf{0.88 s（52\% 短縮）} & \textbf{2.07 s（54\% 短縮）} \\
\end{longtable}

\textbf{考察} - 水平のまま速く旋回すると、遠心力で ZMP が外側に大きく振られ、横方向（最も狭い方向）で転倒する（中央図の破線）。 - 腕を上げても下げても旋回半径（cos α）が小さくなり、慣性と遠心力が減る。ただし \textbf{下げる方が有利} である。 転倒させるモーメントは「高さ × 水平方向の加速度」なので、荷物を低くするとその分だけ効きが弱まる。 - 最適化は、仰角の下限（−35°、地面とのすき間）まで腕を下げ、旋回の中盤（横向き＝最も弱い方向）で最も低くする解を選んだ。 Sim06 の「荷物を体に引き寄せる」と同じく、\textbf{慣性と転倒モーメントを減らす動き} であり、人間やクレーン操作員が経験的に行う動作と一致する。

\section{フィジカル AI への示唆}

\begin{enumerate}
\def\labelenumi{\arabic{enumi}.}
\tightlist
\item
  \textbf{座標の選択はモデル化の一部である。} 球面座標は人間にとって直感的で、旋回動作の指令に向く。 一方で特異点を持ち、学習方策の出力をそのまま関節角にすると、極の近くで関節速度が発散する。 \textbf{方策が何を出力するか（関節角か、手先の向きか）は、学習の安定性と実機の安全性に直結する設計判断} である。
\item
  \textbf{3 次元の姿勢は非可換なので、計画は順序の問題になる。} 平面の結果（Sim01〜05）をそのまま 3 次元に拡張するだけでは不十分で、 系列としての計画・学習が必要になる。
\item
  \textbf{定格荷重マップのような「安全な作業範囲の地図」は、シミュレーションで安価に作れる。} 方向ごとに異なる限界（前後と左右）を学習方策の制約として与えれば、実機で転倒させずに学習を進められる。
\end{enumerate}

\section{限界と次の課題}

\begin{itemize}
\tightlist
\item
  腕は 1 リンク（肘なし）。肘を加えると手先位置と腕の向きを独立に選べ、冗長性を姿勢制御に使える。
\item
  地上モデルは台車が静止している前提（転倒・滑りの開始条件のみ）。
\item
  宇宙での 3 次元の最適ループ（与えた目標姿勢に最短で到達する腕の動作系列）は、まだ求めていない。 非可換性があるため、これは学習・計画の良い課題になる。
\end{itemize}
